\documentclass[11pt,fontset=fandol]{ctexart}
\usepackage[a4paper,margin=23mm]{geometry}
\usepackage{amsmath,amssymb,mathtools,bm,mathrsfs,esint}
\usepackage{graphicx,longtable,booktabs,array,calc}
\usepackage[normalem]{ulem}
\usepackage{xcolor,hyperref}
\hypersetup{hidelinks}
\setcounter{secnumdepth}{-1}
\newcounter{none}
\providecommand{\tightlist}{\setlength{\itemsep}{0pt}\setlength{\parskip}{0pt}}
\providecommand{\passthrough}[1]{#1}
\usepackage{amssymb,mathrsfs,esint,bm,newunicodechar}
\allowdisplaybreaks
\setlength{\emergencystretch}{3em}
\providecommand{\llangle}{\langle\!\langle}
\providecommand{\rrangle}{\rangle\!\rangle}
\newunicodechar{∞}{\ensuremath{\infty}}
\newunicodechar{∈}{\ensuremath{\in}}
\newunicodechar{∑}{\ensuremath{\sum}}
\newunicodechar{∫}{\ensuremath{\int}}
\newunicodechar{≤}{\ensuremath{\leq}}
\newunicodechar{≥}{\ensuremath{\geq}}
\newunicodechar{≠}{\ensuremath{\neq}}
\newunicodechar{→}{\ensuremath{\to}}
\newunicodechar{⇒}{\ensuremath{\Rightarrow}}
\newunicodechar{α}{\ensuremath{\alpha}}
\newunicodechar{β}{\ensuremath{\beta}}
\newunicodechar{γ}{\ensuremath{\gamma}}
\newunicodechar{λ}{\ensuremath{\lambda}}
\newunicodechar{φ}{\ensuremath{\varphi}}
\newunicodechar{ε}{\ensuremath{\varepsilon}}
\newunicodechar{σ}{\ensuremath{\sigma}}
\newunicodechar{θ}{\ensuremath{\theta}}
\newunicodechar{δ}{\ensuremath{\delta}}
\newunicodechar{η}{\ensuremath{\eta}}
\newunicodechar{π}{\ensuremath{\pi}}
\newunicodechar{ξ}{\ensuremath{\xi}}
\newunicodechar{Σ}{\ensuremath{\Sigma}}
\newunicodechar{Δ}{\ensuremath{\Delta}}
\newunicodechar{∂}{\ensuremath{\partial}}
\newunicodechar{∀}{\ensuremath{\forall}}
\newunicodechar{∃}{\ensuremath{\exists}}
\newunicodechar{√}{\ensuremath{\surd}}
\newunicodechar{−}{\ensuremath{-}}
\newunicodechar{μ}{\ensuremath{\mu}}
\newunicodechar{ρ}{\ensuremath{\rho}}
\newunicodechar{τ}{\ensuremath{\tau}}
\newunicodechar{ω}{\ensuremath{\omega}}
\newunicodechar{Ω}{\ensuremath{\Omega}}
\newunicodechar{χ}{\ensuremath{\chi}}
\newunicodechar{ψ}{\ensuremath{\psi}}
\newunicodechar{Γ}{\ensuremath{\Gamma}}
\newunicodechar{Φ}{\ensuremath{\Phi}}
\newunicodechar{Λ}{\ensuremath{\Lambda}}
\newunicodechar{ℝ}{\ensuremath{\mathbb{R}}}
\newunicodechar{ℤ}{\ensuremath{\mathbb{Z}}}
\newunicodechar{ℕ}{\ensuremath{\mathbb{N}}}
\newunicodechar{△}{\ensuremath{\triangle}}
\newunicodechar{⊗}{\ensuremath{\otimes}}
\begin{document}
\section{2023年杭州电子科技大学601数学分析考研真题参考答案}\label{ux5e74ux676dux5ddeux7535ux5b50ux79d1ux6280ux5927ux5b66601ux6570ux5b66ux5206ux6790ux8003ux7814ux771fux9898ux53c2ux8003ux7b54ux6848}

来源：\texttt{HDU/601数学分析/2023/2023年杭州电子科技大学601数学分析考研真题参考答案.pdf}。原
PDF 页序：1, 2, 3, 4, 5, 6, 7, 8, 9。

\section{2023年数学分析打印参考答案：OCR核定整理}\label{ux5e74ux6570ux5b66ux5206ux6790ux6253ux5370ux53c2ux8003ux7b54ux6848ocrux6838ux5b9aux6574ux7406}

\begin{quote}
标注``【OCR公式整理】''的受损段由助手依据同题题面与可核算推导恢复，不作为原作者逐字转录；填空第1题保留打印版与手写版的题面差异。
\end{quote}

\section{杭州电子科技大学2023年数学分析真题解析}\label{ux676dux5ddeux7535ux5b50ux79d1ux6280ux5927ux5b662023ux5e74ux6570ux5b66ux5206ux6790ux771fux9898ux89e3ux6790}

\subsection{一.填空题}\label{ux4e00.ux586bux7a7aux9898}

(本大题共 8 小题,每小题 5 分,共 40 分)

\begin{enumerate}
\def\labelenumi{\arabic{enumi}.}
\tightlist
\item
  打印稿当前识别为 \(f(x)=x/\sqrt{1+x^2}\)，\(f_n\) 是 \(n\) 次复合，求
  \(f_{2023}(0)\)。
\end{enumerate}

【原版差异说明】打印OCR题面和原解答均对应函数值，故此版本下答案为
\(0\)。现有手写原稿第1页可靠识别明确询问导数 \(f'_{2023}(0)\)，其答案为
\(1\)。两份原版存在题面差异，不能把它们合并为同一个问题；手写版的推导另见手写整理稿。

\begin{enumerate}
\def\labelenumi{\arabic{enumi}.}
\setcounter{enumi}{1}
\tightlist
\item
  \(\lim _ { n \to \infty } n ( { \sqrt { n ^ { 2 } + 1 } } - n ) =\)
\end{enumerate}

【解答】

\[
\begin{array} { l } { \displaystyle \lim _ { n \to \infty } n ( \sqrt { n ^ { 2 } + 1 } - n ) = \displaystyle \lim _ { n \to \infty } \frac { n ( \sqrt { n ^ { 2 } + 1 } - n ) ( \sqrt { n ^ { 2 } + 1 } + n ) } { \sqrt { n ^ { 2 } + 1 } + n } } \\ { = \displaystyle \lim _ { n \to \infty } \frac { n } { \sqrt { n ^ { 2 } + 1 } + n } = \displaystyle \lim _ { n \to \infty } \frac { 1 } { \sqrt { 1 + \displaystyle \frac { 1 } { n ^ { 2 } } } + 1 } = \frac { 1 } { 2 } } \end{array}
\]

3.设f(x)在{[}a,b{]}上连续,
\(F ( x ) = \int _ { a } ^ { x } f ( t ) ( t - x ) d t\)
，则F(x)在{[}a,b{]}上的 导数 \(F ^ { \prime } ( x ) =\)

\subsection{【解答】}\label{ux89e3ux7b54}

f(x)在{[}a,b{]}上连续,
\(F ( x ) = \int _ { a } ^ { x } f ( t ) ( t - x ) d t = \int _ { a } ^ { x } t f ( t ) d t - x \int _ { a } ^ { x } f ( t ) d t\)
，故
\(F ^ { \prime } ( x ) = x f ( x ) - \int _ { a } ^ { x } f ( t ) d t - x f ( x ) = - \int _ { a } ^ { x } f ( t ) d t\)

\begin{enumerate}
\def\labelenumi{\arabic{enumi}.}
\setcounter{enumi}{3}
\tightlist
\item
  求 \(\int_{-\infty}^{+\infty}e^{-x^2}\,dx\)。
\end{enumerate}

\subsection{【解答】}\label{ux89e3ux7b54-1}

【OCR公式整理】保留原稿通过圆盘积分夹逼正方形积分的两步计算。

（1）设 \(D_R:x^2+y^2\le R^2\)，则 \[
\iint_{D_R}e^{-x^2-y^2}\,dx\,dy
=\int_0^{2\pi}d\theta\int_0^Re^{-r^2}r\,dr
=\pi(1-e^{-R^2}).
\] 第一象限的四分之一圆盘积分为上述结果的 \(1/4\)。

（2）记 \(J_R=\int_{-R}^Re^{-x^2}\,dx\)。由于
\(D_R\subset[-R,R]^2\subset D_{\sqrt2R}\)，有 \[
\pi(1-e^{-R^2})\le J_R^2
=\int_{-R}^R\int_{-R}^Re^{-x^2-y^2}\,dx\,dy
\le\pi(1-e^{-2R^2}).
\] 夹逼得 \(J_R^2\to\pi\)，又 \(J_R>0\)，故所求积分为 \(\sqrt\pi\)。

\begin{enumerate}
\def\labelenumi{\arabic{enumi}.}
\setcounter{enumi}{4}
\tightlist
\item
  函数 \[
  f(x)=\begin{cases}\cos(\pi x/2),&|x|\le1,\\|x-1|,&|x|>1\end{cases}
  \] 的间断点。
\end{enumerate}

【OCR公式整理】各分支在内部均连续。在
\(x=1\)，两个单侧极限和函数值均为零，故连续。在 \(x=-1\)， \[
\lim_{x\to-1^-}|x-1|=2,\qquad
\lim_{x\to-1^+}\cos(\pi x/2)=0.
\] 故唯一间断点为 \(x=-1\)，是跳跃间断点。原识别左极限中的 \(4\)
已由直接代入核正为 \(2\)。

6.级数
\(\sum _ { n = 2 } ^ { \infty } ( 1 + { \frac { 1 } { n } } ) ^ { n ( n + 1 ) } x ^ { n }\)
的收敛半径为\_

【解答】

\(\lim _ { n \to \infty } \sqrt [ n ] { ( 1 + \frac { 1 } { n } ) ^ { n ( n + 1 ) } } = \lim _ { n \to \infty } ( 1 + \frac { 1 } { n } ) ^ { n + 1 } = e\)
故级数
\(\sum _ { n = 2 } ^ { \infty } ( 1 + { \frac { 1 } { n } } ) ^ { n ( n + 1 ) } x ^ { n }\)
的收敛半径为 \(\frac { 1 } { e }\) 7. 曲线 \(L\) 是圆 \(x^2+y^2=r^2\)
在第一象限的部分，求 \(\int_Lxy\,ds\)。

【OCR公式整理】按半径参数 \(r>0\)，设
\(x=r\cos t,y=r\sin t\)，\(0\le t\le\pi/2\)，有 \[
\int_Lxy\,ds=r^3\int_0^{\pi/2}\cos t\sin t\,dt
=\left[\frac{r^3}{2}\sin^2t\right]_0^{\pi/2}=\frac{r^3}{2}.
\]

\begin{enumerate}
\def\labelenumi{\arabic{enumi}.}
\setcounter{enumi}{7}
\tightlist
\item
  \(\int { \frac { 1 } { x ( 1 + x ^ { n } ) } } d x =\)
\end{enumerate}

\subsection{【解答】}\label{ux89e3ux7b54-2}

\[
\int { \frac { 1 } { x ( 1 + x ^ { n } ) } } d x = \int { \frac { x ^ { n - 1 } } { x ^ { n } ( 1 + x ^ { n } ) } } d x = { \frac { 1 } { n } } \int ( { \frac { 1 } { x ^ { n } } } - { \frac { 1 } { 1 + x ^ { n } } } ) d x ^ { n } = { \frac { 1 } { n } } \ln \left| { \frac { x ^ { n } } { 1 + x ^ { n } } } \right| + C .
\]

\subsection{二.计算题(本大题共5 小题，每小题10分，共 50
分)}\label{ux4e8c.ux8ba1ux7b97ux9898ux672cux5927ux9898ux51715-ux5c0fux9898ux6bcfux5c0fux989810ux5206ux5171-50-ux5206}

\begin{enumerate}
\def\labelenumi{\arabic{enumi}.}
\tightlist
\item
  计算沿上半圆从 \((1,0)\) 到 \((-1,0)\) 的积分 \[
  I=\int_L(-2xe^{-x^2}\sin y)\,dx+(e^{-x^2}\cos y+x^6)\,dy.
  \] 【OCR公式整理】前两项为全微分
  \(d(e^{-x^2}\sin y)\)，两端函数值均为零。补上从 \((-1,0)\) 到
  \((1,0)\) 的直线段 \(L_1\)，其 \(\int x^6\,dy=0\)。由 Green 公式， \[
  \int_Lx^6\,dy=\iint_D6x^5\,dx\,dy=0,
  \] 因为上半圆盘 \(D\) 关于 \(y\) 轴对称。故原积分为零。
\end{enumerate}

2.求
\(\lim _ { x \to \infty } \left[ x - x ^ { 2 } \ln \left( 1 + { \frac { 1 } { x } } \right) \right]\)

【解答】

\[
\lim _ { x \to \infty } \left[ x - x ^ { 2 } \ln \left( 1 + { \frac { 1 } { x } } \right) \right] = \lim _ { x \to \infty } x ^ { 2 } \left[ { \frac { 1 } { x } } - \ln \left( 1 + { \frac { 1 } { x } } \right) \right] = \lim _ { x \to \infty } { \frac { x ^ { 2 } } { 2 x ^ { 2 } } } = { \frac { 1 } { 2 } } .
\]

3.求 \(\sum _ { n = 1 } ^ { \infty } { \frac { 2 n } { 3 ^ { n } } }\)
的和.

\subsection{【解答】}\label{ux89e3ux7b54-3}

幂级数的收敛半径为
\(\lim _ { n \to \infty } { \frac { n } { n + 1 } } = 1\)
，且其在x=±1处通项都是无穷大量，
故其在x=±1处发散，故幂级数的收敛域为(-1,1).

在(-1,1)内，

\[
\sum _ { n = 1 } ^ { \infty } n x ^ { n - 1 } = \sum _ { n = 1 } ^ { \infty } ( x ^ { n } ) ^ { \prime } = ( \sum _ { n = 1 } ^ { \infty } x ^ { n } ) ^ { \prime } = ( \sum _ { n = 0 } ^ { \infty } x ^ { n } ) ^ { \prime } = ( { \frac { 1 } { 1 - x } } ) ^ { \prime } = { \frac { 1 } { ( 1 - x ) ^ { 2 } } }
\]

\(\frac { 1 } { 3 } \in ( - 1 , 1 )\) ，故

\[
\sum _ { n = 1 } ^ { \infty } \frac { 2 n } { 3 ^ { n } } = \frac { 2 } { 3 } \sum _ { n = 1 } ^ { \infty } n ( \frac { 1 } { 3 } ) ^ { n - 1 } = \frac { 2 } { 3 } \frac { 1 } { ( 1 - \frac { 1 } { 3 } ) ^ { 2 } } = \frac { 3 } { 2 } .
\]

\begin{enumerate}
\def\labelenumi{\arabic{enumi}.}
\setcounter{enumi}{3}
\tightlist
\item
  计算 \(\oint_\Gamma x^2y\,dx+y^2\,dy+z\,dz\)，其中 \(\Gamma\) 是
  \(x^2+y^2=1\) 与 \(x+z=1\) 的交线，从 \(z\) 轴正向看为逆时针。
\end{enumerate}

【OCR公式整理】\(y^2\,dy+z\,dz=d(y^3/3+z^2/2)\) 在闭曲线上积分为零。取
\(x=\cos t,y=\sin t,z=1-\cos t\)（\(0\le t\le2\pi\)），则 \[
\oint_\Gamma x^2y\,dx
=-\int_0^{2\pi}\cos^2t\sin^2t\,dt
=-\frac18\int_0^{2\pi}(1-\cos4t)\,dt=-\frac\pi4.
\]

\begin{enumerate}
\def\labelenumi{\arabic{enumi}.}
\setcounter{enumi}{4}
\tightlist
\item
  计算 \(\iiint_\Omega x^2\sqrt{x^2+y^2}\,dx\,dy\,dz\)，其中 \(\Omega\)
  是 \(z=\sqrt{x^2+y^2}\) 与 \(z=x^2+y^2\) 围成的区域。
\end{enumerate}

【OCR公式整理】第二个曲面是抛物面。区域关于交换 \(x,y\) 对称，故原积分为
\[
\frac12\iiint_\Omega(x^2+y^2)^{3/2}\,dV.
\] 保留原参考答案按 \(z\) 分层的思路：\(0\le z\le1\)
时，\(z\le r\le\sqrt z\)。因此 \[
\begin{aligned}
I&=\frac12\int_0^1dz\int_z^{\sqrt z}r^4\,dr\int_0^{2\pi}d\theta\\
&=\frac\pi5\int_0^1(z^{5/2}-z^5)\,dz
=\frac\pi5\left(\frac27-\frac16\right)=\frac\pi{42}.
\end{aligned}
\]

\subsection{三.证明题(本大题共4小题，每小题10分，共40
分)}\label{ux4e09.ux8bc1ux660eux9898ux672cux5927ux9898ux51714ux5c0fux9898ux6bcfux5c0fux989810ux5206ux517140-ux5206}

\begin{enumerate}
\def\labelenumi{\arabic{enumi}.}
\tightlist
\item
  证明 \(\sum_{n=1}^\infty(-1)^n(x^2+n)/n^2\) 在任意有限区间 \([a,b]\)
  一致收敛，在任何实数点均不绝对收敛。
\end{enumerate}

【OCR公式整理】记 \(M=\max\{|a|,|b|\}\)，则 \[
0\le\frac{x^2+n}{n^2}\le\frac{M^2+n}{n^2}\to0
\] 在 \([a,b]\) 上一致成立，且对 \(n\) 单调递减。又
\(\left|\sum_{k=1}^n(-1)^k\right|\le1\)，故由 Dirichlet 判别法一致收敛。
另一方面，对每个实数 \(x\)， \[
\sum_{n=1}^\infty\left|(-1)^n\frac{x^2+n}{n^2}\right|
\ge\sum_{n=1}^\infty\frac1n=\infty.
\] 原识别中绝对值符号遗漏，已按所证命题恢复。

\begin{enumerate}
\def\labelenumi{\arabic{enumi}.}
\setcounter{enumi}{1}
\tightlist
\item
  设 \(f\) 是 \([-1,1]\) 上二阶可导奇函数，\(f(1)=1\)。证明存在
  \(\xi\in(0,1)\)，使 \(f'(\xi)=1\)；并存在 \(\eta\in(-1,1)\)，使
  \(f''(\eta)+f'(\eta)=1\)。
\end{enumerate}

【OCR公式整理】奇性给出 \(f(0)=0\)。Lagrange 中值定理给出
\(f'(\xi)=1\)。因 \(f'\) 为偶函数，\(f'(-\xi)=1\)。 令 \[
F(x)=e^x(f'(x)-1).
\] 则 \(F(\xi)=F(-\xi)=0\)。对 \(F\) 使用 Rolle 定理，存在
\(\eta\in(-\xi,\xi)\)，使 \[
0=F'(\eta)=e^\eta[f''(\eta)+f'(\eta)-1].
\] 所证结论成立。辅助函数中遗漏的 \(f'\)
撇号由这一步求导及端点条件唯一核定。

\begin{enumerate}
\def\labelenumi{\arabic{enumi}.}
\setcounter{enumi}{2}
\tightlist
\item
  设 \(f\) 在 \([0,1]\) 连续，\(D=\{x,y\ge0:x+y\le1\}\)。证明
  \(\iint_Df(1-y)f(x)\,dx\,dy=\frac12[\int_0^1f]^2\)。
\end{enumerate}

【OCR公式整理】保留原解答先换元、后分部积分的思路，令
\(H(x)=\int_x^1f(t)\,dt\)，则 \[
\begin{aligned}
\iint_Df(1-y)f(x)\,dx\,dy
&=\int_0^1 f(x)\left[\int_0^{1-x}f(1-y)\,dy\right]dx\\
&=\int_0^1f(x)H(x)\,dx
=-\frac12[H(x)^2]_0^1
=\frac12\left[\int_0^1f(t)\,dt\right]^2.
\end{aligned}
\]

\begin{enumerate}
\def\labelenumi{\arabic{enumi}.}
\setcounter{enumi}{3}
\tightlist
\item
  已知 \(\Delta\) 表示 Laplace 算符，即 \[
  \Delta u=\frac{\partial^2u}{\partial x^2}+
  \frac{\partial^2u}{\partial y^2}+
  \frac{\partial^2u}{\partial z^2}.
  \] 证明： \[
  \Delta(uv)=(\Delta u)v+2\nabla u\cdot\nabla v+u(\Delta v).
  \] 这里 \(\nabla u\) 表示函数 \(u\) 在点 \((x,y,z)\)
  处的梯度，``\(\cdot\)''表示点乘。
\end{enumerate}

【解答；OCR整理补全】原识别文本中的偏导符号及展开末尾损坏。以下依据本题已识别的目标恒等式，按逐项乘积法则恢复相应公式；本说明由助手添加。

\[
\begin{aligned}
\Delta(uv)
&=\frac{\partial^2(uv)}{\partial x^2}
+\frac{\partial^2(uv)}{\partial y^2}
+\frac{\partial^2(uv)}{\partial z^2}\\
&=(u_{xx}v+2u_xv_x+uv_{xx})
 +(u_{yy}v+2u_yv_y+uv_{yy})
 +(u_{zz}v+2u_zv_z+uv_{zz})\\
&=(u_{xx}+u_{yy}+u_{zz})v
 +2(u_xv_x+u_yv_y+u_zv_z)
 +u(v_{xx}+v_{yy}+v_{zz})\\
&=(\Delta u)v+2\nabla u\cdot\nabla v+u(\Delta v).
\end{aligned}
\]

\subsection{四.综合题(本大题共2小题，共30
分)}\label{ux56db.ux7efcux5408ux9898ux672cux5927ux9898ux51712ux5c0fux9898ux517130-ux5206}

1.（10分）求 \(x+2y=1\) 与 \(x^2+2y^2+z^2=1\) 的交线上距原点最近的点。

【OCR公式整理】保留打印稿消元求最值的思路。约束给出 \[
d=x^2+y^2+z^2=1-y^2,\qquad x=1-2y,\qquad
z^2=4y-6y^2\ge0.
\] 因此 \(0\le y\le2/3\)，\(d\) 的最小值在 \(y=2/3\) 取得，此时
\(x=-1/3,z=0\)，最近点为 \((-1/3,2/3,0)\)。

\begin{enumerate}
\def\labelenumi{\arabic{enumi}.}
\setcounter{enumi}{1}
\tightlist
\item
  设 \(f(x)=\int_1^\infty\cos(xt)/[t(1+t^2)]\,dt\)。证明积分在全实轴关于
  \(x\) 一致收敛，且 \(f\) 一致连续。
\end{enumerate}

【OCR公式整理】保留打印稿的可积控制及积分号下求导思路。对 \(t\ge1\)， \[
\left|\frac{\cos(xt)}{t(1+t^2)}\right|\le\frac1{1+t^2}.
\] 控制函数可积，故原积分关于全体 \(x\) 一致收敛。又 \[
\left|\frac{\partial}{\partial x}
\frac{\cos(xt)}{t(1+t^2)}\right|
=\left|\frac{-\sin(xt)}{1+t^2}\right|
\le\frac1{1+t^2}.
\] 导数积分也一致收敛，故可在积分号下求导： \[
f'(x)=-\int_1^\infty\frac{\sin(xt)}{1+t^2}\,dt,\qquad
|f'(x)|\le\int_1^\infty\frac{dt}{1+t^2}=\frac\pi4.
\]
由中值定理，\(|f(x)-f(y)|\le(\pi/4)|x-y|\)，故全实轴上一致连续。原末尾杂入的分数与游离积分符号已据这一完整估计清理。
\end{document}
